\chapter{Slide Quiz Questions \& Detailed Step-by-Step Code Traces}

\noindent
This section contains all sample quiz questions presented in the official lecture slides (Slides 86 to 105), complete with code analysis and exact step-by-step traces.

\begin{quizbox}[Quiz 1: Inline vs. Internal CSS (Slide 86)]
\textbf{Code}:
\begin{lstlisting}[language=HTML]
<p style="color: red;">Hello CSS!</p>
\end{lstlisting}
\textbf{Question}: Which type of CSS is used in the code above?\\
\textbf{Answer}: \textbf{Inline CSS}\\
\textbf{Explanation}: The \texttt{style} attribute is applied directly inside the HTML opening tag.
\end{quizbox}

\begin{quizbox}[Quiz 2: Internal CSS Location (Slide 87)]
\textbf{Question}: Where should the \texttt{<style>} tag be placed in an HTML document?\\
\textbf{Answer}: \textbf{Inside <head>}\\
\textbf{Explanation}: Internal stylesheets belong in the document head metadata.
\end{quizbox}

\begin{quizbox}[Quiz 3: External CSS Linking (Slide 88)]
\textbf{Code}: \texttt{<link rel="stylesheet" href="style.css">}\\
\textbf{Question}: What does this line do?\\
\textbf{Answer}: \textbf{Imports a CSS file into HTML}
\end{quizbox}

\begin{quizbox}[Quiz 4: CSS Selectors (Slide 89)]
\textbf{Code}:
\begin{lstlisting}[language=HTML]
<p id="demo">Hello World</p>
#demo { color: green; }
\end{lstlisting}
\textbf{Question}: What selector is used here?\\
\textbf{Answer}: \textbf{ID selector} (Identified by the \texttt{\#} prefix).
\end{quizbox}

\begin{quizbox}[Quiz 5: Box Model Outer Spacing (Slide 90)]
\textbf{Code}:
\begin{lstlisting}[language=HTML]
div { margin: 20px; padding: 10px; border: 2px solid black; }
\end{lstlisting}
\textbf{Question}: What is the total outer spacing between the element and its neighbor?\\
\textbf{Answer}: \textbf{20px (margin)}\\
\textbf{Explanation}: Margin defines outer spacing; padding is inner spacing.
\end{quizbox}

\begin{quizbox}[Quiz 6: Background Image Cover (Slide 91)]
\textbf{Code}: \texttt{background-size: cover;}\\
\textbf{Question}: What does \texttt{background-size: cover;} do?\\
\textbf{Answer}: \textbf{Stretches and scales the image to fill the screen/element}.
\end{quizbox}

\begin{quizbox}[Quiz 7: Font Properties (Slide 92)]
\textbf{Code}: \texttt{p \{ font-family: Arial; font-size: 20px; text-align: center; font-weight: bold; \}}\\
\textbf{Question}: What will the paragraph text look like?\\
\textbf{Answer}: \textbf{20px centered bold Arial}.
\end{quizbox}

\begin{quizbox}[Quiz 8: Table Row Hover Effect (Slide 93)]
\textbf{Code}: \texttt{table tr:hover \{ background-color: yellow; \}}\\
\textbf{Question}: What happens when the mouse moves over a table row?\\
\textbf{Answer}: \textbf{The row background turns yellow}.
\end{quizbox}

\begin{quizbox}[Quiz 9: Styling Form Controls (Slide 94)]
\textbf{Code}: \texttt{input, select, textarea \{ border: 2px solid blue; padding: 5px; \}}\\
\textbf{Question}: What will this CSS apply to?\\
\textbf{Answer}: \textbf{All form controls} (\texttt{input}, \texttt{select}, and \texttt{textarea}).
\end{quizbox}

\begin{quizbox}[Quiz 10: Div and Span Cascading (Slide 95)]
\textbf{Code}:
\begin{lstlisting}[language=HTML]
<div style="color: blue;">
  This is a <span style="color: red;">highlighted</span> word.
</div>
\end{lstlisting}
\textbf{Question}: What will be the color of the text ``highlighted''?\\
\textbf{Answer}: \textbf{Red}\\
\textbf{Explanation}: The inline style on \texttt{<span>} directly overrides the inherited color from \texttt{<div>}.
\end{quizbox}

\begin{quizbox}[Quiz 11: Variables and Scope (Slide 96)]
\textbf{Code}:
\begin{lstlisting}[language=JavaScript]
function test() {
  var x = 10;
  if (true) {
    let x = 20;
    console.log(x);
  }
  console.log(x);
}
test();
\end{lstlisting}
\textbf{Question}: What will be the output?\\
\textbf{Answer}: \textbf{20 10}\\
\textbf{Explanation}: \texttt{let x = 20} is scoped strictly to the \texttt{if} block. Outside the block, \texttt{var x = 10} is intact.
\end{quizbox}

\begin{quizbox}[Quiz 12: Using const with Arrays (Slide 97)]
\textbf{Code}:
\begin{lstlisting}[language=JavaScript]
const arr = [1, 2, 3];
arr.push(4);
console.log(arr);
\end{lstlisting}
\textbf{Question}: What is the output?\\
\textbf{Answer}: \textbf{[1, 2, 3, 4]}\\
\textbf{Explanation}: Mutating the contents of a \texttt{const} array is valid; only reassigning the variable reference is forbidden.
\end{quizbox}

\begin{quizbox}[Quiz 13: External vs. Internal Script Execution (Slide 98)]
\textbf{Code}:
\begin{lstlisting}[language=HTML]
<script src="script.js"></script> <!-- contains console.log("External") -->
<script>
  console.log("Internal");
</script>
\end{lstlisting}
\textbf{Question}: What will be printed?\\
\textbf{Answer}: \textbf{``External'' then ``Internal''} (Executed sequentially in order of appearance).
\end{quizbox}

\begin{quizbox}[Quiz 14: Switch Statement Fall-Through (Slide 99)]
\textbf{Code}:
\begin{lstlisting}[language=JavaScript]
let a = 2;
switch(a) {
  case 1: console.log("One");
  case 2: console.log("Two");
  case 3: console.log("Three");
  default: console.log("Default");
}
\end{lstlisting}
\textbf{Question}: What will be the output?\\
\textbf{Answer}: \textbf{Two Three Default}\\
\textbf{Explanation}: Without \texttt{break;} statements, execution falls through \texttt{case 2}, \texttt{case 3}, and \texttt{default}.
\end{quizbox}

\begin{quizbox}[Quiz 15: Loop Logic with Break (Slide 100)]
\textbf{Code}:
\begin{lstlisting}[language=JavaScript]
let i = 0;
while (i < 3) {
  console.log(i);
  i++;
  if (i === 2) break;
}
\end{lstlisting}
\textbf{Question}: What will be printed?\\
\textbf{Answer}: \textbf{0 1}\\
\textbf{Explanation}: Iteration 1 prints 0, increments to 1. Iteration 2 prints 1, increments to 2 $\to$ breaks immediately.
\end{quizbox}

\begin{quizbox}[Quiz 16: Popup Boxes Execution Order (Slide 101)]
\textbf{Code}:
\begin{lstlisting}[language=JavaScript]
let name = prompt("Enter your name:");
if (confirm("Are you sure?")) {
  alert("Welcome, " + name);
} else {
  alert("Action cancelled");
}
\end{lstlisting}
\textbf{Question}: Which popup boxes appear in order?\\
\textbf{Answer}: \textbf{Prompt $\to$ Confirm $\to$ Alert}.
\end{quizbox}

\begin{quizbox}[Quiz 17: For Loop with Continue (Slide 102)]
\textbf{Code}:
\begin{lstlisting}[language=JavaScript]
for (let i = 1; i <= 5; i++) {
  if (i === 3) continue;
  console.log(i);
}
\end{lstlisting}
\textbf{Question}: What numbers will be printed?\\
\textbf{Answer}: \textbf{1 2 4 5} (Number 3 is skipped).
\end{quizbox}

\begin{quizbox}[Quiz 18: Object Property Access \& Unbound this (Slide 103)]
\textbf{Code}:
\begin{lstlisting}[language=JavaScript]
let student = {
  name: "Alice",
  age: 21,
  details: function() {
    return this.name + " is " + this.age + " years old.";
  }
};
let f = student.details;
console.log(f());
\end{lstlisting}
\textbf{Question}: What will be the output?\\
\textbf{Answer}: \textbf{``undefined is undefined years old.''}\\
\textbf{Explanation}: When \texttt{student.details} is stored in \texttt{f} and called unbound, \texttt{this} loses context and refers to the global object/window.
\end{quizbox}

\begin{quizbox}[Quiz 19: do...while Loop Guaranteed Execution (Slide 104)]
\textbf{Code}:
\begin{lstlisting}[language=JavaScript]
let count = 5;
do {
  console.log(count);
  count++;
} while (count < 5);
\end{lstlisting}
\textbf{Question}: What will be the output?\\
\textbf{Answer}: \textbf{5}\\
\textbf{Explanation}: \texttt{do...while} always runs the code block at least once before testing the condition ($6 < 5$ is false).
\end{quizbox}

\begin{quizbox}[Quiz 20: Nested if...else Logic (Slide 105)]
\textbf{Code}:
\begin{lstlisting}[language=JavaScript]
let x = 10, y = 20, z = 30;
if (x > y) {
  if (x > z) console.log("x");
  else console.log("z");
} else {
  if (y > z) console.log("y");
  else console.log("z");
}
\end{lstlisting}
\textbf{Question}: What will be printed?\\
\textbf{Answer}: \textbf{z}\\
\textbf{Explanation}: $x > y$ ($10 > 20$) is false $\to$ jumps to \texttt{else}. $y > z$ ($20 > 30$) is false $\to$ logs ``z''.
\end{quizbox}
